\subsection{Сила натяжения нити}
%%%%%%%%%%%%%%%%%%%%%%%%%%%%%%%%%%%%%%%%%%%%%%%%%%%%
%%%%%%%%%%%%%%%%%%%%%%%%%%%%%%%%%%%%%%%%%%%%%%%%%%%%

 \z{Чтобы передвинуть брусок массой $m=0.6$ кг из состояния покоя вверх вдоль наклонной плоскости на расстояние $l=0.5$ м за время $t=3$ c, надо тянуть нить с постоянной силой $F$. На сколько необходимо увеличить силу $F$, чтобы осуществить такое же перемещение бруска за вдвое меньшее время?}
 %
{$F_1-F=\frac{6lm}{t^2}=0.2\,\text{Н}$}
%%%%%%%%%%%%%%%%%%%%%%%%%%%%%%%%%%%%%%%%%%%%%%%%%%%% 
 \z{Шарик массой $m = 100$ г подвешен на двух одинаковых нитях длиной $a = 1$ м каждая так, что точки подвеса нитей расположены на одной горизонтали. Расстояние между точками подвеса нитей равно $b = 1$ м. Найдите силу натяжения $T$ правой нити сразу после пережигания левой нити.}
 %
{$T=mg\sqrt{1-\left(\frac{b}{2a}\right)^2}\approx 0.87\,\text{Н}$}
%%%%%%%%%%%%%%%%%%%%%%%%%%%%%%%%%%%%%%%%%%%%%%%%%%%% 
 \z[{\begin{wrapfigure}[8]{r}{0.2\textwidth}%
		\vspace{-10pt}%
		\begin{flushright}%
 			\includegraphics[width=0.2\textwidth]{4.5.3}%
		\end{flushright}%
	\end{wrapfigure}%
	}]%
	{Определите ускорение а грузов и силу натяжения $Т$ нитей в системе, изображенной на рисунке. Массы грузов $m_1  = 1$ кг, $m_2  = 2$ кг, $m_3= 3$ кг. Массами нитей и блоков, а также трением пренебречь. Нити нерастяжимы.}
%
{$a=\frac{m_2+m_3-m_1}{m_1+m_2+m_3}g=6.5\,\frac{\text{м}}{\text{с}^2};$~$T_1=\frac{2m_1 (m_2+m_3)g}{m_1+m_2+m_3}=16\,\text{Н};$~$T_3=\frac{2m_1 m_3 g}{m_1+m_2+m_3}=10\, \text{Н}$}
%%%%%%%%%%%%%%%%%%%%%%%%%%%%%%%%%%%%%%%%%%%%%%%%%%%%
 \z {На одном конце веревки, переброшенной через невесомый блок, находится груз массой $m$, а на другом -- обезьянка массой $2m$. Она поднимается вертикально вверх с ускорением $a_\text{отн}  = g$ относительно веревки. Каково ее ускорение $а$ относительно земли?}
 %
{$a=0$}
%%%%%%%%%%%%%%%%%%%%%%%%%%%%%%%%%%%%%%%%%%%%%%%%%%%% 
 \z {Три резиновых шнура связывают вместе и медленно растягивают в разные стороны (рис. слева). В некоторый момент длины всех трёх шнуров оказываются равны $L_1  = 20$ см. Затем шнуры растягивают под другими углами (рис. справа). В этом случае равенство длин шнуров наступает при длине $L_2  = 30$ см каждого из них.  Известна начальная длина самого длинного шнура в недеформированном состоянии: $l =15$ см. Найдите длины двух других шнуров и отношение жесткостей шнуров.
 		\begin{center}%
 			\includegraphics[width=0.8\textwidth]{4.5.5}%
		\end{center}%
 }
%
{$k_1:k_2:k_3=1:1:1.78;$~$l_1=l_2=L_1-\frac{k_3}{k_1}(L_1-l)\approx 11\,\text{см}$}
%%%%%%%%%%%%%%%%%%%%%%%%%%%%%%%%%%%%%%%%%%%%%%%%%%%%